\documentclass[12pt,a4paper]{article}

\usepackage[utf8]{inputenc}
\usepackage[T1]{fontenc}
\usepackage[brazil]{babel}
\usepackage[a4paper,margin=2.cm]{geometry}
\usepackage{enumitem}
\setlist[enumerate]{itemsep=0.45em,topsep=0.35em}

\begin{document}

\begin{center}
  {\Large\bfseries PROVA SIMULADA -- NPOR}\par
  \vspace{0.4cm}
  Nome: \rule{9cm}{0.4pt}\par\vspace{0.3cm}
  Data: \rule{4cm}{0.4pt}\hfill Nota: \rule{3cm}{0.4pt}
\end{center}

\noindent\textbf{Duração sugerida: 90 minutos.} A prova oficial pode ter outro tempo e outros conteúdos; confirme as instruções da sua Organização Militar. Este simulado tem 13 questões objetivas e uma redação.

\medskip
\noindent\textbf{Instruções:} marque uma alternativa por questão. Na redação, escreva um texto dissertativo-argumentativo, com letra legível, respeitando o tema e o espaço disponível.

\subsection*{Português}
Leia o texto para responder às questões 1 e 2.

\begin{quote}
Em períodos de preparação intensa, é comum querer estudar muitas horas de uma só vez. No entanto, uma rotina curta e constante permite perceber dúvidas, corrigi-las e voltar ao conteúdo com mais clareza. O planejamento não substitui o esforço; ajuda a direcioná-lo.
\end{quote}

\begin{enumerate}
  \item Qual é a ideia principal do texto?
  \begin{enumerate}[label=\alph*)]
    \item Estudar por muitas horas seguidas é a única forma de aprender.
    \item O planejamento elimina a necessidade de esforço.
    \item Uma rotina constante ajuda a direcionar o estudo e identificar dúvidas.
    \item Toda pessoa deve estudar exatamente o mesmo número de horas.
  \end{enumerate}

  \item No trecho ``ajuda a direcioná-lo'', o pronome ``lo'' retoma:
  \begin{enumerate}[label=\alph*)]
    \item o período.
    \item o esforço.
    \item o conteúdo.
    \item o planejamento.
  \end{enumerate}

  \item Assinale a frase com concordância verbal adequada.
  \begin{enumerate}[label=\alph*)]
    \item Fazem dois meses que comecei a revisão.
    \item Houveram mudanças no cronograma.
    \item Existem assuntos que exigem mais prática.
    \item Os candidatos chegou cedo, obrigatoriamente.
  \end{enumerate}

  \item Assinale a alternativa em que o uso da crase está correto.
  \begin{enumerate}[label=\alph*)]
    \item Entreguei o documento à secretaria.
    \item Começou à estudar pela manhã.
    \item Voltou à pé para casa.
    \item Pedi ajuda à ele.
  \end{enumerate}
\end{enumerate}

\subsection*{Matemática e raciocínio lógico}
\begin{enumerate}[resume]
  \item Um material custa R\$ 160,00 e recebe desconto de 15\%. Qual é o preço final?
  \begin{enumerate}[label=\alph*)]
    \item R\$ 124,00
    \item R\$ 136,00
    \item R\$ 145,00
    \item R\$ 150,00
  \end{enumerate}

  \item Se 4 cadernos custam R\$ 30,00, quanto custam 10 cadernos pelo mesmo preço unitário?
  \begin{enumerate}[label=\alph*)]
    \item R\$ 60,00
    \item R\$ 70,00
    \item R\$ 75,00
    \item R\$ 80,00
  \end{enumerate}

  \item Resolva a equação $3x-7=11$.
  \begin{enumerate}[label=\alph*)]
    \item $x=4$
    \item $x=6$
    \item $x=7$
    \item $x=8$
  \end{enumerate}

  \item Um retângulo mede 8 cm por 5 cm. Qual é sua área?
  \begin{enumerate}[label=\alph*)]
    \item 13 cm$^2$
    \item 26 cm$^2$
    \item 40 cm$^2$
    \item 80 cm$^2$
  \end{enumerate}

  \item Complete a sequência: 3, 7, 11, 15, \ldots
  \begin{enumerate}[label=\alph*)]
    \item 17
    \item 18
    \item 19
    \item 20
  \end{enumerate}
\end{enumerate}

\subsection*{Conhecimentos gerais e institucionais}
\begin{enumerate}[resume]
  \item Em que ano foi proclamada a Independência do Brasil?
  \begin{enumerate}[label=\alph*)]
    \item 1789
    \item 1822
    \item 1889
    \item 1930
  \end{enumerate}

  \item A Força Expedicionária Brasileira (FEB) atuou em qual conflito?
  \begin{enumerate}[label=\alph*)]
    \item Primeira Guerra Mundial.
    \item Guerra do Paraguai.
    \item Segunda Guerra Mundial.
    \item Guerra de Canudos.
  \end{enumerate}

  \item Qual alternativa apresenta as cinco regiões oficiais do Brasil?
  \begin{enumerate}[label=\alph*)]
    \item Norte, Nordeste, Centro-Oeste, Sudeste e Sul.
    \item Amazônica, Sertaneja, Central, Litorânea e Sul.
    \item Norte, Leste, Centro, Sudeste e Oeste.
    \item Amazônia, Nordeste, Pantanal, Sudeste e Pampas.
  \end{enumerate}

  \item Qual alternativa lista apenas símbolos nacionais previstos na legislação brasileira?
  \begin{enumerate}[label=\alph*)]
    \item Bandeira, Hino, Armas e Selo Nacionais.
    \item Bandeira, Constituição, brasão estadual e Hino.
    \item Hino, lema do Exército, Bandeira e uniforme.
    \item Selo, mapa do Brasil, Bandeira e continência.
  \end{enumerate}
\end{enumerate}

\newpage
\section*{Redação}
\noindent\textbf{Tema:} A importância da disciplina e do planejamento na formação de jovens comprometidos com a sociedade.

\medskip
\noindent Escreva um texto dissertativo-argumentativo, apresentando seu ponto de vista, argumentos que o sustentem e uma conclusão. Sugestão: de 15 a 25 linhas. Não é necessário escrever um título.

\vspace{0.4cm}
\noindent\rule{\textwidth}{0.3pt}\par\vspace{0.45cm}
\noindent\rule{\textwidth}{0.3pt}\par\vspace{0.5cm}
\noindent\rule{\textwidth}{0.3pt}\par\vspace{0.5cm}
\noindent\rule{\textwidth}{0.3pt}\par\vspace{0.5cm}
\noindent\rule{\textwidth}{0.3pt}\par\vspace{0.5cm}
\noindent\rule{\textwidth}{0.3pt}\par\vspace{0.5cm}
\noindent\rule{\textwidth}{0.3pt}\par\vspace{0.5cm}
\noindent\rule{\textwidth}{0.3pt}\par\vspace{0.5cm}
\noindent\rule{\textwidth}{0.3pt}\par\vspace{0.5cm}
\noindent\rule{\textwidth}{0.3pt}\par\vspace{0.5cm}
\noindent\rule{\textwidth}{0.3pt}\par\vspace{0.5cm}
\noindent\rule{\textwidth}{0.3pt}\par\vspace{0.5cm}
\noindent\rule{\textwidth}{0.3pt}\par\vspace{0.5cm}
\noindent\rule{\textwidth}{0.3pt}\par\vspace{0.5cm}
\noindent\rule{\textwidth}{0.3pt}\par\vspace{0.5cm}
\noindent\rule{\textwidth}{0.3pt}\par\vspace{0.5cm}
\noindent\rule{\textwidth}{0.3pt}\par\vspace{0.5cm}
\noindent\rule{\textwidth}{0.3pt}\par\vspace{0.5cm}
\noindent\rule{\textwidth}{0.3pt}\par\vspace{0.5cm}
\noindent\rule{\textwidth}{0.3pt}\par\vspace{0.5cm}
\noindent\rule{\textwidth}{0.3pt}\par\vspace{0.5cm}
\noindent\rule{\textwidth}{0.3pt}\par

\end{document}